\section{Motivations}
\label{subsec:gear-design-objectives}

The design of GEAR was motivated by the heterogeneity of existing AI-based multi-agent frameworks and by the implicit nature of their configuration spaces. Although these frameworks provide capabilities for defining agents, tasks, models, tools, memory mechanisms, and orchestration structures, they expose these capabilities through different abstractions, terminologies, and implementation mechanisms. Moreover, the available configuration options and the constraints governing their combinations are generally distributed across documentation, APIs, configuration files, and source code. Based on these observations GEAR was designed to address three limitations encountered when configuring and experimenting with agentic systems: the implicit representation of their configuration spaces, the fragmentation of their concepts across frameworks, and the difficulty of reproducing equivalent system designs across heterogeneous execution platforms. These limitations were translated into the following design objectives.

\paragraph{\textbf{Explicitly represent the configuration space of
agentic systems.}}
Existing frameworks expose numerous configuration decisions concerning
agents, tasks, models, tools, memory mechanisms, orchestration structures,
and execution policies. These decisions are generally distributed across
source code, API calls, configuration files, and documentation. Their
dependencies and compatibility constraints consequently remain largely
implicit. GEAR should provide an explicit representation of these variation
points and their constraints, allowing configurations to be systematically
constructed, inspected, and validated before execution.

\paragraph{\textbf{Separate system specifications from
framework-specific implementations.}}
Agentic frameworks frequently provide conceptually similar capabilities
through different abstractions, terminology, and APIs. As a result, a system
architecture designed for one framework cannot generally be reused in
another without manually reinterpreting and reimplementing its configuration.
GEAR should therefore provide a framework-independent specification model
that captures the architecture and configurable properties of an agentic
system without embedding the implementation details of a particular target
framework.

\paragraph{\textbf{Support automated and reproducible instantiation
across frameworks.}}
A framework-independent specification is useful only if it can be
transformed into concrete implementations. GEAR should therefore validate a
configuration, translate its concepts into the abstractions of a selected
framework, and deterministically generate the corresponding executable
artifacts. Given the same GEAR specification, target framework, connector
version, and execution environment, the transformation process should
produce the same implementation structure.

These objectives led to a separation between two complementary concerns.
The first concerns the specification of the system and its variability,
including its agents, reusable orchestration modules, workflows, and
configuration constraints. The second concerns the processing of this
specification through validation, framework-specific translation,
instantiation, and execution. The construction of these two parts is
described in the following subsections.


